\section{Introduction}
\label{introduction}


Penetration Testing (short {\em pentesting}) is a methodology for
assessing network security, by generating and executing possible
attacks exploiting known vulnerabilities of operating systems and
applications (e.g., \cite{ArcGra04}). Doing so automatically allows
for regular and systematic testing without a prohibitive amount of
human labor, and makes pentesting more accessible to non-experts. A
key question is how to automatically generate the attacks.


A natural way to address this issue is as an {\em attack planning}
problem. This is known in the AI Planning community as the ``Cyber
Security'' domain \cite{BodGohHaiHar05}. Independently (though
considerably later), the approach was put forward also by Core
Security \cite{LucSarRic10}, a company from the
pentesting industry. In that form, attack planning is very technical,
addressing the low-level system configuration details that are
relevant to vulnerabilities. %
Herein, we are concerned exclusively with this setting. We consider
regular automatic pentesting as done in Core Security's ``Core Insight
Enterprise'' tool. We will use the term ``attack planning'' in that
sense.

Lucangeli et al.\ \shortcite{LucSarRic10} encode attack planning into
PDDL, and use off-the-shelf planners. This already is useful---in
fact, it is currently employed commercially in Core Insight
Enterprise, using a variant of Metric-FF
\cite{hoffmann:jair-03}. However, the approach is limited by its
inability to handle uncertainty. The pentesting tool cannot be
up-to-date regarding all the details of the configuration of every
machine in the network, maintained by individual users.

Core Insight Enterprise currently addresses this by extensive use of
{\em scanning} methods as a pre-process to planning, which then
considers only {\em exploits}, \ie, hacking actions modifying the
system state. The drawbacks of this are that (a) this pre-process
incurs significant costs in terms of running time and network traffic,
and (b) even so, since scans are not perfect, a residual uncertainty
remains (Metric-FF is run based on the configuration that appears to
be most likely). Prior work \cite{SarRicLuc11} has addressed (b) by
associating each exploit with a success probability. This is unable to
model dependencies between the exploits, % 
and it still requires extensive scanning (to obtain realistic success
probabilities) so does not solve (a). Herein, we provide the first
solution able to address both (a) and (b), intelligently mixing scans
with exploits like a real hacker would. The basic insight is that
penetration testing can be naturally modeled in terms of solving a
POMDP.


We encode the incomplete knowledge as an uncertainty of state, thus
modeling the possible network configurations in terms of a probability
distribution. Scans and exploits are deterministic in that their
outcome depends only on the state they are executed in. Negative
rewards encode the cost (the duration) of scans and exploits; positive
rewards encode the value of targets attained. The model incorporates
firewalls, detrimental side-effects of exploits (crashing programs or
entire machines), %
and dependencies between exploits relying on similar vulnerabilities.


POMDP solvers fail to scale to large networks. This is not
surprising---even the input model grows exponentially in the number of
machines. We show how to address this based on exploiting network
structure. We view networks as graphs whose vertices are
fully-connected subnetworks, and whose arcs encode the connections
between these, filtered by firewalls.  We decompose this graph into
biconnected components. We approximate the attacks on these components
by combining attacks on individual subnetworks. We approximate the
latter by combining attacks on individual machines. The approximations
are conservative, \ie, they never over-estimate the value of the
policy returned. %
Attacks on individual machines are modeled and solved as POMDPs, and
the solutions are propagated back up. We evaluate this approach based
on the test suite of Core Insight Enterprise, showing that, compared
to a global POMDP model, it vastly improves runtime at a small cost in
attack quality.

We next discuss some preliminaries. We then describe our POMDP model,
our decomposition algorithm, and our experimental findings, before
concluding the paper.
